\section*{Chapitre \ref{ch:b2:amines} --- Amines et composés azotés}

\begin{solution}{exo:b2:amines:1}
Primaire~; secondaire~; tertiaire~; sel d'ammonium quaternaire~; primaire (\omterm{def:b2:huckel:aromatic}{aromatique}).
\end{solution}

\begin{solution}{exo:b2:amines:2}
Éthanamide $<$ aniline $<$ ammoniac $<$ méthylamine ($\mathrm pK_a$ des acides conjugués
$-0.6$, 4.60, 9.26, 10.66).
\end{solution}

\begin{solution}{exo:b2:amines:3}
\textit{N}-méthylcyclohexanimine, \ce{C6H10=N-CH3}~: amine et cétone en milieu faiblement acide
(pH 4--5), en éliminant l'eau formée (desséchant ou distillation azéotropique).
\end{solution}

\begin{solution}{exo:b2:amines:4}
Bromobenzène~; benzonitrile~; phénol~; 4-hydroxyazobenzène \ce{C6H5-N=N-C6H4-OH} (couplage en para du
groupe \ce{O-} du phénolate).
\end{solution}

\begin{solution}{exo:b2:amines:5}
$\mathrm pK_a = 9.80$~: à pH 7, le rapport $[\ce{(CH3)3NH+}]/[\ce{(CH3)3N}] = 10^{2.8} \approx 630$~: la
forme protonée prédomine~; le jus de citron (pH voisin de 2) la protone pratiquement en totalité.
Agiter une solution organique d'une amine avec un acide aqueux dilué fait passer l'amine, sous
forme d'ion ammonium, dans l'eau~; une base la libère de nouveau.
\end{solution}

\begin{solution}{exo:b2:amines:6}
1-(cyclohex-1-én-1-yl)pyrrolidine. Après l'addition et la perte d'eau, l'azote porte une charge
positive (un ion iminium) et aucun hydrogène à perdre~; c'est le carbone voisin qui perd un
proton, ce qui forme la liaison \ce{C=C}.
\end{solution}

\begin{solution}{exo:b2:amines:7}
Le benzaldéhyde avec la propan-2-amine, ou la propanone avec la benzylamine, chaque fois avec
\ce{NaBH3CN} vers pH 6.
\end{solution}

\begin{solution}{exo:b2:amines:8}
La propiophénone (1-phénylpropan-1-one), \ce{C6H5COCH2CH3}~: une seule addition donne l'anion
d'\omterm{def:b2:amines:imine}{imine}, hydrolysé en cétone.
\end{solution}

\begin{solution}{exo:b2:amines:9}
Phtalimide de potassium + 1-bromohexane~: \textit{N}-hexylphtalimide~; l'hydrazine (ou une hydrolyse)
libère ensuite l'hexan-1-amine. L'azote du phtalimide alkylé n'a pas de doublet disponible (il est
délocalisé sur deux carbonyles) et plus aucune liaison \ce{N-H}~: il ne peut pas être alkylé une seconde fois.
\end{solution}

\begin{solution}{exo:b2:amines:10}
L'eau de brome brome l'aniline aux trois positions ortho et para du groupe amino~: le
2,4,6-tribromoaniline se forme. La \omterm{def:b2:amines:diazonium}{diazotation}, puis \ce{H3PO2}, remplace le groupe amino par H~: les trois
bromes sont désormais en 1,3,5 les uns par rapport aux autres.
\end{solution}

\begin{solution}{exo:b2:amines:11}
4-nitroaniline + \ce{NaNO2} + \ce{HCl} à 0--\qty{5}{\degreeCelsius}~: chlorure de
4-nitrobenzènediazonium. Couplage avec le phénol en solution faiblement basique (phénolate, fortement
activé) à la position para de \ce{O-}~: le 4-[(4-nitrophényl)diazényl]phénol, \ce{O2N-C6H4-N=N-C6H4-OH},
un colorant orange.
\end{solution}

\begin{solution}{exo:b2:amines:12}
Le groupe triméthylammonium est un groupe partant médiocre et volumineux, et il rend les hydrogènes $\beta$
du côté \ce{CH3} plus acides et plus accessibles~; la base arrache un hydrogène du carbone le moins
substitué~: le but-1-ène, l'alcène le moins substitué (produit de Hofmann).
\end{solution}

\begin{solution}{pb:b2:amines:1}
\textbf{1.} $\mathrm{{}^-O_3S{-}C_6H_4{-}NH_3^+}$~: le groupe sulfonique, acide, a protoné l'amine.
\textbf{2.} L'amphion est peu soluble dans l'eau~; le carbonate déprotone le groupe ammonium, ce qui
donne le sulfanilate de sodium, soluble.
\textbf{3.} \ce{NaNO2 + HCl -> HNO2 + NaCl}.
\textbf{4.}
\[
  \mathrm{{}^-O_3S{-}C_6H_4{-}NH_2 + HNO_2 + H^+ \to {}^-O_3S{-}C_6H_4{-}N_2^+ + 2\,H_2O} .
\]
\textbf{5.} Le sel de diazonium se décompose à chaud (en phénol et \ce{N2})~; l'acide nitreux
se décompose aussi.
\textbf{6.} $5.00/173.19 = \qty{28.9}{mmol}$~: \qty{28.9}{mmol} de \ce{NaNO2}, soit \qty{1.99}{g}.
\textbf{7.} L'acide nitreux en excès réagirait avec le partenaire de couplage~; on le détecte avec du papier
iodo-amidonné et on le détruit par l'urée ou l'acide sulfamique.
\textbf{8.} Son cycle est fortement activé par le groupe diméthylamino, un donneur puissant.
\textbf{9.} En para de \ce{N(CH3)2}~: orientation ortho/para, et les positions ortho sont encombrées par les
groupes méthyle.
\textbf{10.} $\mathrm{{}^-O_3S{-}C_6H_4{-}N{=}N{-}C_6H_4{-}N(CH_3)_2}$.
\textbf{11.} En milieu fortement acide, le groupe diméthylamino est protoné~: \ce{-NH(CH3)2+} désactive le cycle
et le couplage s'arrête.
\textbf{12.} Le produit se forme sous sa forme acide protonée, en partie soluble~; en milieu basique, le
sulfonate de sodium, moins soluble dans la solution salée, précipite.
\textbf{13.} Sa charge positive est délocalisée sur le cycle et l'azote terminal est un centre électrophile
médiocre~; seuls les cycles activés par \ce{-O-}, \ce{-OH} ou \ce{-NR2} réagissent.
\textbf{14.} Jaune orangé en milieu basique, rouge en milieu acide.
\textbf{15.} Sur un azote du groupe azo (la forme protonée est stabilisée par conjugaison avec le
groupe diméthylamino).
\textbf{16.} Vers $\mathrm pK_a \pm 1$~: de 3.1 à 4.4 environ, autour de $\mathrm pK_a = 3.47$.
\textbf{17.} Deux cycles et le groupe azo forment un long système conjugué~: l'écart entre le plus haut niveau
$\pi$ occupé et le plus bas vacant est faible (\cref{prop:b2:huckel:gap}) et l'absorption
se situe dans le visible.
\textbf{18.} La protonation modifie la répartition des charges dans le système $\pi$ et réduit encore l'écart~:
l'absorption se déplace vers les grandes longueurs d'onde, la couleur passe du jaune au rouge.
\textbf{19.} \qty{173.19}{g/mol}.
\textbf{20.} \qty{327.33}{g/mol}.
\textbf{21.} \qty{28.87}{mmol}.
\textbf{22.} $0.02887 \times 327.33 = \qty{9.45}{g}$.
\textbf{23.} La décomposition du sel de diazonium, un couplage incomplet, la solubilité du produit dans
les eaux mères et les eaux de lavage.
\textbf{24.} $0.80 \times 9.45 = \boldsymbol{\qty{7.56}{g}}$ d'hélianthine.
\end{solution}
